\begin{abstract}
Let $f:Y\to B$ be a proper holomorphic map with connected fibres
from a smooth compact complex threefold with $b_1(Y)=b_2(Y)=0$,
and suppose that its smooth general fibre is a complex two-torus.
We prove that every reduced fibre support has nonnegative Euler
number, with arbitrary component multiplicities and arbitrary
support singularities. Consequently $B\simeq\mathbb P^1$ and
$c_3(Y)\in\{0,2\}$. Finite-character support cohomology computes
the Euler number of an irreducible support with rational or ruled
normalization resolution; local derivative and conductor
comparisons identify the zero-Euler normalizations. For reducible
supports, the actual Picard restrictions exclude elliptic--Enoki
interfaces, and a local vanishing sphere excludes the remaining
all-Enoki case. The conductor Euler formula then gives both
nonnegativity and the complete equality criterion. The supporting
occurrence constructions, multiplicity calculations, contractions
and comparison examples are included with their precise maps.
\end{abstract}

\paragraph{The proof's scope.}
The irreducible argument uses the local comparisons of Paper~II
for the single support in question, together with the
finite-character mixed-Hodge calculation proved here. Its use
does not require the other supports to be irreducible or to have
normal crossings. For a singular irreducible support of Euler
number zero, the actual normalization is smooth minimal elliptic
ruled or smooth minimal Hopf. For a reducible zero-Euler support,
the actual normalization components are smooth minimal Hopf or
non-K\"ahler elliptic surfaces and their reduced conductor maps
are finite unramified covers of disjoint smooth elliptic curves.

\paragraph{Known inputs and the additional analytic form.}
The value $e(Y)=2$ for the supplied completed torus family
is stated in the original manuscript
\cite[Theorem on p.~4; Theorem~7.22, p.~60]{AC26}.
Engel computes this value from the two zero-dimensional
toric strata in \cite[Proposition~7.1]{Engel2026}.
The relative cusp construction is also stated in
\cite[Section~2, relative cusp theorem]{DGKX2026v1,DGKX2026v2}.
These existing constructions and Euler computations retain
their credit.  The nonnegativity and equality conclusions
of \ref{cdp:st:C.2.G.1} are proved for every family satisfying
its displayed hypotheses by the local comparisons below.
The classification in
\cite[Section~1 and Theorem~3.1]{KamenovaTori} assumes
a semistable K\"ahler total space with trivial canonical
bundle; each use of that classification retains those
hypotheses.
The surface-classification facts are cited at
\cite[Proposition~2.1 and Theorem~2.2, pp.~842--843]{Peters2021}.
The logarithmic support complex uses
\cite[Section~3, Theorem~1 and Corollary~2; Section~4, equation~(3)]{Steenbrink1999}.
That source works with a complete algebraic ambient space.
The compact-projective-strata construction required for the
analytic neighbourhood here is proved explicitly in
\ref{cdp:st:C.2.MH.1}; the source is not cited as already
containing that broader form. The class-VII foliation and
negative-definite curve inputs are cited at the exact loci
opened in Favre and Dloussky--Oeljeklaus.

The cyclic symplectic torus orders used in
\ref{cdp:st:C.2.UN.24} and \ref{cdp:st:C.2.UN.32}
are classical \cite[Lemmas~3.2--3.3, pp.~30--31]{Fujiki1988Tori}.
The four cyclic quotient configurations are listed in
\cite[table on p.~7659; Theorem~27, p.~7663]{Roulleau2019Kummer},
with their credit to Nikulin and Bertin retained.
For an algebraic-dimension-zero K3 surface, the cyclic
torus quotient associated with a foliation is already proved in
\cite[Section~10, pp.~29--30]{Brunella2002Foliations}.
The stabilizer action on the saturated limiting-field space
in \ref{cdp:st:C.2.UN.16} has an eigenline.  Its tangent line
descends through the finite quotient on the unramified locus,
then extends by saturation on the K3 resolution, giving the
foliation to which that result applies.  The elliptic-cover
case in \ref{cdp:st:C.2.UN.32} uses the explicit translation
quotient and principal-bundle characteristic-class calculation
to construct its torus.  The occurrence statements retain the
specified quotient component through the actual vertical
modifications.

\paragraph{Reading the verification record.}
Historical checks, current re-derivations, partial current
reviews and source openings are recorded separately. The long
occurrence catalogue has its own per-statement review record;
the assembly of the main theorem supplies no blanket check for
those constructions.
